\documentclass[dvipdfmx,autodetect-engine]{jsarticle}
\usepackage{amsmath,amssymb,bm}
\usepackage{enumitem}
\usepackage[dvipdfmx]{hyperref}

\title{Kano and Takai (2011) 注釈付き日本語読解ノート}
\author{読解対象：Yutaka Kano and Keiji Takai\\
\small Journal of Multivariate Analysis 102 (2011), 1241--1255}
\date{}

\newcommand{\E}{\mathbb{E}}
\newcommand{\Prb}{\mathbb{P}}
\newcommand{\Cov}{\operatorname{Cov}}
\newcommand{\Var}{\operatorname{Var}}
\newcommand{\indep}{\mathrel{\perp\!\!\!\perp}}
\newcommand{\note}[1]{\par\smallskip\noindent\textbf{注釈：}#1\par\smallskip}
\newcommand{\connect}[1]{\par\smallskip\noindent\textbf{欠測IVメモへの接続：}#1\par\smallskip}
\newcommand{\cue}[1]{\par\smallskip\noindent\textbf{原文の短い手がかり：}#1\par\smallskip}

\begin{document}
\maketitle

\section*{書誌情報と読解方針}

本ノートは，Kano and Takai による ``Analysis of NMAR missing data without specifying missing-data mechanisms in a linear latent variate model'' を，日本語で順に読むための注釈付き資料である。対象論文は \textit{Journal of Multivariate Analysis} 102(2011), 1241--1255 に掲載された。

本ノートは原論文の逐語的な全文翻訳ではなく，論文全体の議論を冒頭から末尾まで追えるようにした詳細な日本語パラフレーズである。英語原文は短い確認句に限定し，式・仮定・推定法・数値例・文献上の位置づけを中心に再構成する。

\subsection*{用語対応表}

\begin{center}
\begin{tabular}{p{0.38\linewidth}p{0.52\linewidth}}
NMAR & missing not at random，非無作為欠測 \\
latent variate model & 潜在変数モデル，線形潜在変量モデル \\
shared-parameter model & 共有パラメータモデル \\
complete-case analysis (CCA) & 完全ケース分析，リストワイズ削除後の分析 \\
multi-sample analysis (MSA) & 多母集団分析，多標本SEM \\
pattern-mixture model & パターン混合モデル \\
selection model & 選択モデル \\
factorial invariance & 因子不変性 \\
asymptotic robustness & 漸近ロバスト性 \\
\end{tabular}
\end{center}

\connect{本プロジェクトの \texttt{manuscript/main.tex} は，欠測が潜在因子または観測アウトカムに依存する場合の識別を扱う。Kano and Takai (2011) は，欠測が潜在変数に依存するNMARを，欠測機構を明示的に推定せず，SEMの多母集団分析として扱う点で重要である。}

\section{Abstract}

\cue{latent variables / not missing at random / without specifying mechanisms}

要旨は，潜在変数モデルで欠測機構が潜在変数に依存するのは自然だが，その場合はMARではなくNMARになる，という問題意識から始まる。潜在変数は測定誤差を除いた概念であり，応用研究者が関心を持つ構成概念を表すため，欠測がそこに依存すると考えるのは自然である。

論文の提案は，欠測機構の具体的な関数形を推定に使わず，線形潜在変量モデルの全パラメータに対して一致性と漸近正規性を持つ推定法を与えることである。この方法は，平均構造あり・なしの多母集団分析として実装できるため，SEMソフトウェアで扱いやすい。

さらに，古典的には避けるべきとされがちな完全ケース分析も，このモデルでは因子負荷量や誤差分散など一部の重要パラメータに対して一致推定量を与えることが示される。

\note{この論文の「欠測機構を指定しない」は，欠測が無視可能という意味ではない。欠測は潜在変数に依存するのでNMARだが，条件付き分布の構造を使って，欠測パターン別のSEMとして推定する。}

\section{1. Introduction}

導入部は，不完全データを扱う二つの古典的分解から始まる。選択モデルは同時分布を $\Prb(R\mid Y;\tau)f(Y\mid\theta)$ と分解し，パターン混合モデルは $f(Y\mid R;\theta)\Prb(R\mid\tau)$ と分解する。ここで $R$ は欠測指標ベクトル，$Y$ は観測変数ベクトル，$\theta$ は関心パラメータ，$\tau$ は欠測機構側のパラメータである。

欠測パターンは $r^{(\ell)}$ で表され，完全ケースは $r^{(1)}=(1,\ldots,1)'$ である。各パターンでは観測される変数だけを選ぶ選択行列 $D_{r^{(\ell)}}$ が使われ，実際の観測ベクトルは $D_RY$ と書かれる。

MARは，欠測パターンの確率が未観測成分ではなく観測成分だけに依存する条件である。MARなら，$Y_{\mathrm{obs}}$ の尤度に基づく直接最尤法，SEMでいうfull-information maximum likelihoodが $\theta$ の一致推定量を与える。NMARでは，この直接尤度は一般に一致性を保証しない。

潜在変数モデルでは，選択モデルとパターン混合モデルに加えて，共有パラメータモデルが使われてきた。連続潜在変数 $z$ または潜在クラス $c$ を共有し，$Y$ と $R$ が $z$ または $c$ を条件づけると独立になるという考え方である。

既存の共有パラメータモデルでは，欠測機構の具体的な関数形を明示して推定することが多い。しかし適切な欠測機構を見つけるのは難しい。本論文は，共有パラメータ的な条件付き独立を仮定しつつ，推定では欠測機構の関数形を使わない。代わりに，パターン混合モデルの発想で，欠測パターン別の条件付き分布を分析する。

論文の構成は，2節で潜在変数に依存する欠測機構がNMARであることを示し，3節でSEMにおける新しい推定法を提案し，4節で漸近理論を扱い，5節で数値例を示し，6節で文献上の位置づけと結論を述べる。

\section{2. Missing mechanism depending on latent variables}

この節では，確認的因子分析モデル
\[
Y=\mu+\Lambda(\lambda)z+e,\qquad \E(z)=0,\quad \E(e)=0
\]
を考える。共通因子 $z$ と独自因子または誤差 $e$ は無相関で，分散は $\Var(z)=\Phi$，$\Var(e)=\Psi$ である。観測変数の分散は
\[
\Var(Y)=\Lambda(\lambda)\Phi\Lambda(\lambda)'+\Psi=\Sigma(\theta)
\]
と表される。

\begin{description}[leftmargin=3.2em]
\item[A1] $Y$ と $R$ は，潜在変数 $z$ を条件づけると独立である。
\end{description}

A1は，欠測が観測変数 $Y$ の測定誤差成分ではなく，共通因子 $z$ によって説明されるという仮定である。つまり
\[
\Prb(R=r\mid Y,z)=\Prb(R=r\mid z)=h(r\mid z;\tau)
\]
と書ける。$h$ が $z$ に依存しなければMCARであるが，論文が扱うのは $h$ が $z$ に依存する場合である。

正規性を仮定すると，$\Prb(R=r^{(\ell)}\mid Y)$ は $z\mid Y$ の条件付き正規密度で $h(r^{(\ell)}\mid z)$ を平均したものになる。一方，MARで必要なのは $\Prb(R=r^{(\ell)}\mid D_{r^{(\ell)}}Y)$ で同じ確率が表せることである。因子分析モデルでは，$z\mid Y$ は通常すべての $Y$ の成分に依存するため，観測成分だけを条件づけた分布とは異なる。したがってMAR条件は一般に成立しない。

\note{ここでの論点は重要である。欠測が「真の能力」「潜在態度」「景気因子」などに依存するなら，観測指標の一部だけを条件づけても欠測確率を説明しきれない。これは自然なNMARである。}

\section{3. Analysis of incomplete data in SEM}

3節は，より一般の線形潜在変量モデル
\[
Y=\mu+\sum_{g=1}^G\Lambda_g(\lambda)z_g+e
\]
を扱う。$z_g$ は外生潜在ベクトル，$e$ は誤差ベクトルで，相互に無相関である。分散構造は
\[
\Var(Y)=\sum_{g=1}^G\Lambda_g(\lambda)\Phi_g\Lambda_g(\lambda)'+\Psi=\Sigma(\theta)
\]
である。SEMの多くはこの枠組みに含まれる。

MAR下で使われる直接尤度は，欠測パターンごとに観測成分の正規密度を使うが，$\theta$ と $\mu$ は全パターンで共通とする。この直接尤度は，この論文のNMAR状況では正当化されない。

\begin{description}[leftmargin=3.2em]
\item[A2] $z_1,\ldots,z_G,e$ は独立であり，$R$ と $Y$，または $R$ と $e$ は $(z_1,\ldots,z_G)$ を条件づけると独立である。
\item[A3] 欠測確率は潜在変数群を通じて分解され，
\[
\Prb(R=r\mid Y,z_1,\ldots,z_G)=\Prb(R=r\mid z_1,\ldots,z_G)=\prod_{g=1}^Gh_g(r\mid z_g;\tau,\theta)
\]
と書ける。
\end{description}

A3の第一等号はA2の帰結である。A3は，欠測機構が観測変数 $Y$ と無関係であるとは言っていない。$Y$ と $R$ は潜在変数を共有して依存する。潜在変数は $Y$ と $R$ の交絡因子として働く。

A2とA3の下で，欠測パターン $R=r^{(\ell)}$ を条件づけても，潜在変数 $z_g$ の独立性が保たれる。各 $z_g\mid R=r^{(\ell)}$ は，元の密度を $h_g(r^{(\ell)}\mid z_g)$ で傾けた分布になる。この性質から，
\[
Y\mid R=r^{(\ell)}
\]
も，同じ因子負荷行列 $\Lambda_g(\lambda)$ と同じ誤差分散 $\Psi$ を持つ線形潜在変量モデルに従う。ただし，潜在平均 $\nu_g^{(\ell)}$ と潜在分散 $\Phi_g^{(\ell)}$ はパターンごとに変わる。

この再生性が論文の核心である。欠測機構が観測変数に直接依存する場合には，一般にこの性質は成立しない。

\subsection*{完全ケース分析}

完全ケース $R=r^{(1)}$ だけを見ると，$Y\mid R=r^{(1)}$ は同じ $\Lambda(\lambda)$ と $\Psi$ を持つ因子モデルに従う。したがって完全ケース数が十分大きく，関心が因子負荷量や誤差分散にあるなら，正規尤度をモーメント法の道具として使うCCAは一致推定を与える。

ただし，完全ケースの潜在分散 $\Phi^{(1)}$ は母集団の $\Phi$ ではない。したがってCCAは因子共分散や因子平均を一般には正しく推定できない。

\subsection*{提案法：パターン別多母集団分析}

全データを使うために，論文は欠測パターンごとに条件付きSEMを立てる。尤度は
\[
L_{\mathrm{PTM}}(\theta,\nu\mid Y)
\]
として表され，パターンごとに観測成分 $D_{r^{(\ell)}}Y$ の正規密度を使う。共通にするのは因子負荷量 $\lambda$ と誤差分散 $\Psi$ であり，潜在平均 $\nu_g^{(\ell)}$ と潜在分散 $\Phi_g^{(\ell)}$ はパターンごとに自由にする。

この方法はパターン混合モデルに似ているが，実装上は平均構造付きSEMの多母集団分析である。因子不変性モデルと同じく，群間で因子負荷量と誤差分散を等置し，潜在平均・潜在分散は群ごとに変える。

潜在平均と一般平均には不定性があるため，識別制約が必要である。論文は，自然な制約
\[
\sum_{\ell=1}^L\Prb(R=r^{(\ell)})\nu_g^{(\ell)}=0
\]
または，完全ケース群の潜在平均を0に置く単純な制約を議論する。

\begin{description}[leftmargin=3.2em]
\item[A4] 各欠測パターンにおいて，観測成分が潜在平均を識別できるだけの情報を持つ。
\end{description}

A4は，欠測パターン内で観測されている指標が少なすぎると潜在平均を推定できない，という当然の条件を形式化する。たとえば，ある潜在因子の全指標が同時に欠測しているパターンでは，その因子の平均や分散をそのパターンから推定できない。

母集団の潜在分散 $\Phi_g$ は，全分散分解
\[
\Phi_g=\E\{\Var(z_g\mid R)\}+\Var\{\E(z_g\mid R)\}
\]
により，各パターンの $\Phi_g^{(\ell)}$ と $\nu_g^{(\ell)}$ から復元できる。これは平均構造付きMSAを使う大きな利点である。

MAR直接尤度との違いは明確である。MAR直接尤度はすべてのパターンで平均・潜在分散まで等しいと制約する。提案法は，$\lambda$ と $\Psi$ だけを等置し，欠測パターンごとの潜在平均・潜在分散を許す。したがって，MAR直接尤度は提案法に強い等値制約を追加した特殊ケースとして見られる。

\note{提案法は欠測機構 $h_g$ の関数形を推定しない。その代わり，欠測パターンによって潜在分布が変わることを許す。これは「選択で傾いた潜在因子分布」をパターンごとの母集団として扱う発想である。}

\connect{\texttt{main.tex} では潜在因子 $F$ が欠測に関わる場合，単純な $R\indep V\mid \bm Y$ は破れると整理している。Kano and Takaiは，その場合でも線形潜在変数モデルの構造を使えば，因子負荷量や誤差分散は推定可能であることを示す。}

\section{4. Asymptotic properties of estimators}

4節は，欠測パターン別の標本平均と標本分散が，ランダムなパターン別標本サイズにもかかわらず通常の多母集団分析と同じ漸近理論を持つことを示す。

母集団分布はかなり一般に取り，$Y\mid R$ の4次モーメントが有限であるとする。パターン $\ell$ で観測される平均と分散を
\[
\mu_\ell=\E(D_{r^{(\ell)}}Y\mid R=r^{(\ell)}),\qquad
\Sigma_\ell=\Var(D_{r^{(\ell)}}Y\mid R=r^{(\ell)})
\]
と定義する。

標本では，各観測がどの欠測パターンに入るかを示す指示変数 $M_i^{(\ell)}$ を使う。パターン別サンプルサイズ $n_\ell$ はランダムだが，標本平均 $\bar Y_\ell$ と標本分散 $S_\ell$ はそれぞれ $\mu_\ell$ と $\Sigma_\ell$ に一致する。

中心極限定理により，
\[
\sqrt{n_\ell}\{(\bar Y_\ell-\mu_\ell)',\operatorname{vec}(S_\ell-\Sigma_\ell)'\}'
\]
は漸近正規となり，異なるパターン間では漸近的に独立になる。したがって，一つの不完全データ標本を欠測パターン別に分けた解析は，通常の多母集団SEMと同じ漸近理論で扱える。

提案尤度の正規理論形式は，パターンごとの標本平均・標本分散を使った目的関数として書ける。平均が利用できない場合には，分散構造だけの尤度を使い，$\lambda$ と $\Psi$ の一致推定だけを狙う。

推定量の漸近分散は，通常のM推定量やSEMの多母集団分析の理論と同じ sandwich 型で表される。正規分布が真であれば簡略化された分散式が得られる。

さらに重要なのは，線形潜在変量モデルの独立性構造と誤差分布の仮定により，完全ケース分析の正規理論標準誤差や適合度検定が，非正規母集団に対しても漸近ロバスト性を持つ場合がある，という点である。欠測パターンで選択された標本であっても，条件付きモデルの構造が保たれるため，Browne and Shapiro 型のロバスト性が適用できる。

\note{この節の技術的貢献は，パターン別サンプルサイズがランダムでも，多母集団SEMとして扱えることを明確にした点である。実装上は通常のSEMソフトを使えるが，理論上はランダム群サイズの扱いが必要になる。}

\section{5. Numerical example}

数値例では，NMAR欠測を持つ人工データに対して，CCA，平均構造ありMSA，平均構造なしMSA，MAR直接ML，完全データ分析を比較する。目的は，提案法がFisher一致性を持つこと，MAR直接MLがNMARでは必ずしも一致しないこと，そしてパターンごとの潜在平均・分散が大きく異なることを確認することである。

モデルは，9個の観測変数と3個の共通因子を持つ確認的因子分析である。各因子に3つの指標が対応し，因子負荷量の一部は1に固定される。欠測は $Y_2,Y_5,Y_8$ に発生し，それぞれ対応する潜在因子 $z_1,z_2,z_3$ が閾値を超えると欠測する。したがって8つの欠測パターンが生じ，A3を満たす。

表1は，各欠測パターンにおける潜在平均と潜在分散を示す。条件付き潜在平均・分散はパターン間で大きく異なり，条件付き潜在分散は母集団分散をかなり過小評価する。このことは，CCAだけで因子分散を推定すると偏る理由を示している。

表2は因子負荷量推定を比較する。完全データ分析が参照値であり，平均構造ありMSAが最も良好で，MAR直接MLがそれに続く。CCAと平均構造なしMSAはこの標本ではやや劣る。小さい欠測パターンだけの分析は標本サイズ不足のため不安定になる。

ただし，MAR直接MLがこの人工データで見かけ上よく働いても，母集団条件付き平均・分散を直接代入した解析ではFisher一致でないことが確認される。これは，NMARではMAR直接MLが理論的に一致しないという主張と整合する。

表3は，平均構造ありMSAで推定されたパターン別の潜在平均・分散を示す。識別のため完全ケース群の潜在平均を0に固定しているため，表1の真の条件付き平均とは見かけが異なる。しかし，母集団因子分散の復元には問題ない。

表4は因子分散行列 $\Phi$ の推定結果である。CCAは $\Phi$ をうまく推定できない。平均構造ありMSAは，分散分解を使って良好な推定を与える。MAR直接MLも標本上は近い結果を与えるが，Fisher一致ではない。

\note{数値例のメッセージは，因子負荷量と誤差分散だけならCCAも理論的に成立しうるが，因子分散や因子平均まで欲しいなら，欠測パターン別の潜在平均・分散を推定するMSAが必要，ということである。}

\section{6. Discussion and conclusion}

議論は，観測変数には測定誤差があり，潜在変数は研究者の関心概念をより直接に表すため，欠測を潜在変数の関数としてモデル化するのは自然である，という立場から始まる。この仕様は合理的で単純だが，欠測はNMARになり，通常は欠測機構の明示的モデリングが必要になる。

本論文は，この状況に対して四つの方法を検討した。完全ケース分析，平均構造ありMSA，平均構造なしMSA，MAR直接MLである。A2--A4の下では，平均構造ありMSAが全パラメータに対して妥当な一致推定量を与える。実装は通常のSEM多母集団分析なので，欠測機構を含む完全尤度をMonte Carloで評価するより容易である。

提案MSAはMAR直接MLに似ているが，欠測パターン間で潜在平均と潜在分散を等置しない点が違う。その分パラメータ数は多い。また，このMSAはMARだがMCARではない欠測には必ずしも適さない可能性があり，MARと本論文型NMARを経験的に区別する方法や，両方に使えるハイブリッド手法が今後の課題として挙げられる。

方法の弱点は，A4が満たされない場合があること，および観測変数が多く同時に欠測するパターンでは潜在平均や潜在分散を推定できないことである。すべての変数が欠測するパターンは情報を持たないため無視せざるを得ない。SEMソフトによっては，各パターンの標本分散行列の正定値性を要求するため，観測変数数よりパターン内サンプルサイズが大きい必要もある。

階層的因子分析モデルを例に，A3の積構造は多様な欠測機構を表せることが説明される。欠測が上位一般因子だけに依存する場合，下位グループ因子に依存する場合，項目群ごとの欠測が対応する潜在因子に依存する場合などを含む。

文献レビューでは，MAR下のFIML，planned missing design，因子不変性，潜在クラス共有パラメータモデル，未知NMAR機構に対する擬似尤度，完全ケース分析の正当化，ロジット型NMARモデル，補助変数によるMAR近似，非正規母集団でのロバスト性などが整理される。

結論として，適切な欠測機構 $\Prb(R\mid Y)$ または $\Prb(R\mid z)$ が見つかり，識別・推定できるなら選択モデルでMLを行えばよい。しかし多くの場合それは難しい。その場合，本論文型のNMARでは，平均構造ありの部分等値制約MSAが有力な代替となる。完全ケースがほとんどならCCAも選択肢だが，潜在分散のバイアスには注意が必要である。MAR直接MLも使える場合があるが，無条件に推奨するにはさらなる検討が必要である。

\section{Appendix: Vec operators}

付録は，本文の漸近分散式で使われる行列演算を整理する。$\operatorname{vec}(A)$ は行列の列を縦に積んだベクトル，$v(A)$ は対称行列の異なる要素だけを集めたベクトルである。重複行列 $D_p$ は $v(A)$ を $\operatorname{vec}(A)$ に写す線形作用素であり，交換行列 $K_{pq}$ は $\operatorname{vec}(A)$ を $\operatorname{vec}(A')$ に変換する。対称化行列 $N_p=(I_{p^2}+K_{pp})/2$ も導入される。

正規ベクトル $X\sim N_p(0,\Sigma)$ について，
\[
\Var\{\operatorname{vec}(XX')\}=2N_p(\Sigma\otimes\Sigma)
\]
が成り立つ。この式が，正規理論SEMの標本分散行列の漸近分散を簡潔に書くために使われる。

\section{Referencesの位置づけ}

参考文献は，Rubin，Little and Rubin，Schafer などの欠測データ基礎文献，Bollen，Anderson，Browne and Shapiro，Satorra and Bentler などのSEM・線形潜在変数モデル・漸近ロバスト性文献，Muthen and Asparouhov，Lee and Tang，Song and Lee などの潜在変数モデルにおけるNMAR文献に分かれる。d'Haultfoeuille型の欠測IVとは異なるが，欠測が潜在構造を通じて発生するときの識別・推定可能性を示す点で，本プロジェクトに直接関係する。

\section*{全体まとめ}

Kano and Takai (2011) の主張は，欠測が潜在変数に依存するNMARであっても，線形潜在変量モデルでは条件付きモデルの構造が欠測パターンごとに再生されるため，欠測機構を明示的に推定せずにSEMの多母集団分析で推定できる，というものである。共通に保つべきものは因子負荷量と誤差分散であり，潜在平均と潜在分散は欠測パターンごとに変える。完全ケース分析も一部パラメータには有効だが，潜在分散や平均の推定には限界がある。この論文は，欠測IVの識別方程式とは別のルートで，NMAR欠測と潜在変数モデルを結びつける重要な基礎文献である。

\end{document}
